overuse of brackets, 
too many times mentioned 49 papers from 2022 to 2025, 
for too simple things you dont need to explain again a bracket

Column matches (for example, determining that \texttt{zip\_code} in one table means the same as \texttt{postal\_code} in another), entity matches or clusters (groups of records that refer to the same entity), and merged datasets. -> same example twice, you have explained column match, entity match once before this you are repeating?

you again mentioned postal_code zip_code example column match in table 3 -> remove,

if you have explained the term once if you dont have to use brackets or anything to explain that gaian

do-no-harm behavior -> what do you mean by this

approximate key relationship -> explain

efficiency as first-class concerns. -> use something more survey paper worthy language instead of first class concern

you have defined schema, entity matching, schema matching in multiple place which is unnecessary, define everything once when you first mention not any time after that, that looks stupid

too many times you ahve used zip code postal code example, just do once and remove others

first you mentioned like SIND then like sINDs is there any reason, if its correct then keep it otherwise unify it


in figure 2 yDI signal is mentioned, define what that is


phone number email example mentioned more than once, keep only one


earlier work is cited but not covered in the same depth -> I did not see any earlier work before 2022 being cited, check this infor remove if the info is wrng


paiwise match accuray -> define the term


(3)~fairness-aware evaluation (checking whether matching errors are distributed unevenly across demographic groups), and (4)~reusable integration signals (lightweight evidence, such as approximate foreign-key relationships, that can be shared across tasks).-> i didnot understand

TabSketchFM~\cite{11113110} avoids the cost of converting entire tables into text (which can exceed the input limits of language models) by compressing each table into a compact statistical sketch that preserves key similarity signals. A transformer model is trained to read these sketches directly and produce embeddings for efficient retrieval. -> i did not understand

LakeBench~\cite{10.14778/3659437.3659448} provides labeled ground-truth queries paired with workloads that test both retrieval quality and system efficiency, enabling fair comparison of discovery methods. -> i did not understand

that package cleaning know-how into automated, controllable steps. -> i did not understand


two product listings might be considered the same for inventory purposes but different for price comparison---and models this as multi-label classification. -> did not understand

(for example, ``age'' when studying the effect of a drug)---and retrieves them from external sources. It then constructs a causal graph that guides which variables to integrate. -> did not understand


combines embeddings with practical indexings -> did not understand


